\chapter{Recursos e cronograma}
\label{cap:recursos}

\section{Recursos}
\label{sec:recursos}

\subsection{Recursos humanos}

A execução é conduzida pelo autor, estudante de graduação em Engenharia Elétrica,
sob orientação do Prof.\ Dr.\ Leandro dos Santos Coelho, do Departamento de
Engenharia Elétrica da UFPR e do PPGEE/UFPR, responsável pela condução
metodológica nas frentes de modelagem substituta, otimização por metaheurísticas e
regressão simbólica. A coorientação do Prof.\ Dr.\ Bruno Pohlot Ricobom, do
Departamento de Engenharia Elétrica da UFPR, cobre a frente de instrumentação e
medição: a concepção das placas de teste, a operação do escâner de campo próximo e
a interpretação dos resultados da validação experimental complementar descrita na
Seção~\ref{sec:validacao_exp}.

Estima-se dedicação de dez horas semanais ao longo de dois semestres letivos,
com reuniões de acompanhamento quinzenais com o orientador.

\subsection{Recursos materiais e computacionais}

O trabalho não requer aquisição de equipamentos nem licenças comerciais. Toda a
cadeia de ferramentas é livre, conforme relacionado na
Tabela~\ref{tab:ferramentas}.

\begin{table}[htb]
\centering
\caption{Recursos computacionais e de software}
\label{tab:ferramentas}
\footnotesize
\begin{tabular}{@{}p{0.26\linewidth}p{0.30\linewidth}p{0.36\linewidth}@{}}
\toprule
\textbf{Recurso} & \textbf{Especificação} & \textbf{Finalidade} \\
\midrule
Estação de trabalho & computador pessoal, 8\,GB de RAM & desenvolvimento,
treinamento e análise \\[0.3em]
Armazenamento & 60\,GB livres em disco & base bruta (24\,GB), base processada e
resultados \\[0.3em]
Python 3.12 & \texttt{numpy}, \texttt{pandas}, \texttt{scipy} & manipulação
numérica \\[0.3em]
\texttt{scikit-rf} & --- & leitura de arquivos Touchstone e conversão entre
parâmetros \\[0.3em]
\texttt{scikit-learn}, \texttt{LightGBM} & --- & modelos de referência \\[0.3em]
\texttt{PyTorch} & versão para CPU & modelo informado por física \\[0.3em]
\texttt{Optuna} & --- & otimização de hiperparâmetros e dos pesos $\lambda_k$
\\[0.3em]
\texttt{pymoo} & --- & NSGA-II e evolução diferencial multiobjetivo \\[0.3em]
\texttt{SHAP}, \texttt{PySR} & --- & interpretabilidade e regressão simbólica
\\[0.3em]
\texttt{MLflow} & --- & registro de experimentos \\[0.3em]
\texttt{pytest}, \texttt{Hydra} & --- & testes automatizados e configuração
\\[0.3em]
Git & repositório privado & versionamento de código e do texto \\[0.3em]
\LaTeX & modelo UFPR & redação do documento \\
\bottomrule
\end{tabular}
\vspace{0.3em}
\\{\footnotesize FONTE: O autor (2026).}
\end{table}

O dimensionamento computacional foi verificado quanto à viabilidade. O
treinamento de um perceptron multicamadas sobre 985 amostras de oito dimensões,
com saída de 334 pontos, é executável em CPU em poucos minutos. O gargalo real é
a etapa de pré-processamento, que percorre 24\,GB de arquivos Touchstone; medições
preliminares indicam tempo de leitura da ordem de 0{,}1\,s por arquivo, o que
projeta cerca de dois minutos para a base completa em uma passagem única, cujo
resultado é persistido. Como plano de contingência, caso a memória disponível se
mostre insuficiente na etapa de otimização, prevê-se o uso do Google Colab em
sua camada gratuita.

\subsection{Recursos financeiros}

Não há custos previstos para o escopo principal. Os dados são de acesso gratuito
e o software é livre. Um custo eventual, estimado entre R\$\,300 e R\$\,600,
poderia incorrer na fabricação das duas placas de teste da validação experimental
complementar, item opcional que não integra o caminho crítico e cujo
financiamento seria discutido com o laboratório de origem caso venha a ser
executado.

\section{Cronograma}
\label{sec:cronograma}

O trabalho distribui-se por dois semestres letivos: 2026/2, dedicado à
fundamentação, à infraestrutura de dados e ao estabelecimento das referências; e
2027/1, dedicado ao modelo informado por física, à otimização, à
interpretabilidade e à redação final. A divisão respeita a estrutura de TCC I e
TCC II do curso.

\subsection{Marcos de decisão}

Adotam-se quatro marcos com critérios objetivos de continuidade, apresentados na
Tabela~\ref{tab:marcos}. A finalidade é permitir redirecionamento antecipado, em
vez de constatação tardia de inviabilidade.

\begin{table}[htb]
\centering
\caption{Marcos de decisão e planos de contingência}
\label{tab:marcos}
\footnotesize
\begin{tabular}{@{}p{0.06\linewidth}p{0.17\linewidth}p{0.32\linewidth}p{0.34\linewidth}@{}}
\toprule
& \textbf{Prazo} & \textbf{Critério de continuidade} & \textbf{Contingência} \\
\midrule
M1 & set./2026 & unidades confirmadas; pipeline de extração de $Z(f)$ validado
contra invariantes analíticas em toda a base (a amostra preliminar de 40
configurações já é consistente) & rever hipótese de unidades; consultar os
mantenedores da base \\[0.4em]
M2 & nov./2026 & melhor referência atinge $R^2 > 0{,}80$ sobre $|Z|_{\max}$ &
revisar representação do alvo; considerar redução de dimensionalidade da curva
\\[0.4em]
M3 & abr./2027 & modelo informado por física supera a ablação em ao menos uma
fração de treinamento, com significância estatística & reportar resultado
negativo como contribuição; deslocar ênfase para interpretabilidade \\[0.4em]
M4 & jun./2027 & fronteira de Pareto verificada por simulação de referência &
restringir a otimização mono-objetivo sobre $|Z|_{\max}$ \\
\bottomrule
\end{tabular}
\vspace{0.3em}
\\{\footnotesize FONTE: O autor (2026).}
\end{table}

Cabe registrar que o marco M3 não condiciona a validade do trabalho. A hipótese
H1 é falsificável, e sua refutação — a constatação de que a informação física
não produz ganho mensurável neste problema — constitui resultado legítimo e
publicável, desde que o desenho experimental tenha poder estatístico suficiente
para sustentá-la. O plano de contingência associado preserva a entrega, e não
apenas a confirmação da hipótese.

\subsection{Detalhamento por período}

\begin{table}[htb]
\centering
\caption{Cronograma de execução --- TCC I (2026/2)}
\label{tab:crono1}
\footnotesize
\begin{tabular}{@{}p{0.52\linewidth}ccccc@{}}
\toprule
\textbf{Atividade} & \textbf{ago} & \textbf{set} & \textbf{out} & \textbf{nov} &
\textbf{dez} \\
\midrule
Revisão bibliográfica sistemática              & \mes & \mes & \mes &      &      \\
Caracterização e validação da base de dados    & \mes & \mes &      &      &      \\
Pipeline de extração de impedância e testes    &      & \mes & \mes &      &      \\
Análise exploratória dos dados                 &      &      & \mes &      &      \\
Implementação dos modelos de referência        &      &      & \mes & \mes &      \\
Protocolo experimental e registro de execuções &      &      &      & \mes &      \\
Redação dos capítulos 1 a 3                    &      &      & \mes & \mes & \mes \\
Entrega e apresentação de TCC I                &      &      &      &      & \mes \\
\bottomrule
\end{tabular}
\vspace{0.3em}
\\{\footnotesize FONTE: O autor (2026).}
\end{table}

\begin{table}[htb]
\centering
\caption{Cronograma de execução --- TCC II (2027/1)}
\label{tab:crono2}
\footnotesize
\begin{tabular}{@{}p{0.52\linewidth}ccccc@{}}
\toprule
\textbf{Atividade} & \textbf{mar} & \textbf{abr} & \textbf{mai} & \textbf{jun} &
\textbf{jul} \\
\midrule
Formulação e implementação dos termos físicos  & \mes & \mes &      &      &      \\
Calibração dos pesos e curvas de aprendizado   &      & \mes & \mes &      &      \\
Testes de hipótese H1 e H2                     &      &      & \mes &      &      \\
Otimização multiobjetivo do empilhamento       &      &      & \mes & \mes &      \\
Interpretabilidade e teste da hipótese H3      &      &      &      & \mes &      \\
Validação experimental complementar (opcional) &      &      &      & \mes &      \\
Redação dos capítulos de resultados e conclusão&      &      & \mes & \mes & \mes \\
Verificação de integridade e revisão final     &      &      &      &      & \mes \\
Defesa                                         &      &      &      &      & \mes \\
\bottomrule
\end{tabular}
\vspace{0.3em}
\\{\footnotesize FONTE: O autor (2026).}
\end{table}

\section{Resultados fundamentais a serem atingidos}
\label{sec:resultados_esperados}

Ao final do trabalho, esperam-se as seguintes entregas:

\begin{alineas}
\item um modelo substituto informado por física para predição da curva de
impedância de PDNs de seis camadas, com $R^2 > 0{,}90$ sobre $|Z|_{\max}$ em
conjunto de teste isolado e latência de inferência inferior a 1\,ms;

\item uma quantificação do efeito da informação física sobre o erro de
generalização em função do volume de dados de treinamento, com significância
estatística avaliada, respondendo às hipóteses H1 e H2;

\item uma fronteira de Pareto de empilhamentos que equilibram impedância máxima e
custo de fabricação estimado, com as soluções extremas verificadas por simulação
de referência;

\item um conjunto de expressões analíticas interpretáveis relacionando
parâmetros de empilhamento ao risco eletromagnético, acompanhado de faixas de
validade, respondendo à hipótese H3;

\item uma implementação de código aberto, reprodutível e documentada, com testes
automatizados e registro de experimentos\seapendice{, conforme as práticas
descritas no Apêndice~\ref{ap:engenharia}};

\item o texto do TCC segundo as normas da UFPR e, como desdobramento pretendido,
um manuscrito submetido a evento da área — IEEE EMC+SIPI ou DesignCon — bem como
base para ingresso no mestrado do PPGEE/UFPR.
\end{alineas}
